\scanpage{108}
\mspage{53}
canonique
\[
(152)\qquad D_{\unsure{q|d}}\longrightarrow X
\]
\margin{On introduit aussi $D_0=X$, et $D_0\longrightarrow X$ l'identité}
qui se factorise d'ailleurs par $D$, et qui est un morphisme d'immersion local, de codim $q$. \unsure{Sur|Pour} $D_{\unsure{q|d}}$, on a un \unsure{revêtt|revêtement} étale à $q$ feuillets, noté $\overset{\deleted{\unsure{q}}}{\widetilde D_{\unsure{q|d}}}$, \deleted{formé des couples}\inserted{dont la fibre en} un pt $(x,V)$ est \unsure{formée|formé} de l'ens des $\unsure{q|d}$ branches $D_{xi}$ de $D$ en $x$ qui contiennent $V$. Donc les \deleted{\unsure{points|couples}} pts de $\widetilde D_{\unsure{q|d}}$ correspondent aux triplets $(x,\unsure{V},W)$, où $x$ est un pt de $D$, et $V\subset W$ des germes de ss-multiplicités de $X$ en $x$ attachées à $D$, de codims respectives $\unsure{q|d}$ et $1$. On trouve donc un morphisme canonique
\[
\diagram{\begin{tikzcd}[ampersand replacement=\&, row sep=small]
\widetilde D_{\unsure{q|d}}\arrow[r] \& D_1\\
(x,V,W)\arrow[r,mapsto] \& W
\end{tikzcd}}
\]

Plus généralement, si $J$ est un ens. fini ordonné comme ci-dessus, \inserted{muni d'une fonction codim,} on \deleted{\unsure{en}} désigne par $D_J$ la \deleted{\unsure{var.}} multiplicité \deleted{formée}\inserted{dont} les points sont les couples d'un $x\in\deleted{D}\inserted{X}$ \inserted{de codim $d=d(J)$,} et d'un \unsure{plongt|plongement} d'ens. ordonnés, compatible \deleted{\unsure{avec|aux}} fonctions codim
\margin{on suppose la fonction codim strictement décr. cf (151)}
\[
(152)\qquad J\mathrel{\overset{\varphi}{\hookrightarrow}}\Delta_x\simeq\mathcal P(I(x))^\circ
\]
\margin{107}
